% Options for packages loaded elsewhere
\PassOptionsToPackage{unicode}{hyperref}
\PassOptionsToPackage{hyphens}{url}
\PassOptionsToPackage{dvipsnames,svgnames,x11names}{xcolor}
%
\documentclass[
  letterpaper,
  DIV=11,
  numbers=noendperiod]{scrartcl}

\usepackage{amsmath,amssymb}
\usepackage{lmodern}
\usepackage{iftex}
\ifPDFTeX
  \usepackage[T1]{fontenc}
  \usepackage[utf8]{inputenc}
  \usepackage{textcomp} % provide euro and other symbols
\else % if luatex or xetex
  \usepackage{unicode-math}
  \defaultfontfeatures{Scale=MatchLowercase}
  \defaultfontfeatures[\rmfamily]{Ligatures=TeX,Scale=1}
\fi
% Use upquote if available, for straight quotes in verbatim environments
\IfFileExists{upquote.sty}{\usepackage{upquote}}{}
\IfFileExists{microtype.sty}{% use microtype if available
  \usepackage[]{microtype}
  \UseMicrotypeSet[protrusion]{basicmath} % disable protrusion for tt fonts
}{}
\makeatletter
\@ifundefined{KOMAClassName}{% if non-KOMA class
  \IfFileExists{parskip.sty}{%
    \usepackage{parskip}
  }{% else
    \setlength{\parindent}{0pt}
    \setlength{\parskip}{6pt plus 2pt minus 1pt}}
}{% if KOMA class
  \KOMAoptions{parskip=half}}
\makeatother
\usepackage{xcolor}
\setlength{\emergencystretch}{3em} % prevent overfull lines
\setcounter{secnumdepth}{-\maxdimen} % remove section numbering
% Make \paragraph and \subparagraph free-standing
\ifx\paragraph\undefined\else
  \let\oldparagraph\paragraph
  \renewcommand{\paragraph}[1]{\oldparagraph{#1}\mbox{}}
\fi
\ifx\subparagraph\undefined\else
  \let\oldsubparagraph\subparagraph
  \renewcommand{\subparagraph}[1]{\oldsubparagraph{#1}\mbox{}}
\fi

\usepackage{color}
\usepackage{fancyvrb}
\newcommand{\VerbBar}{|}
\newcommand{\VERB}{\Verb[commandchars=\\\{\}]}
\DefineVerbatimEnvironment{Highlighting}{Verbatim}{commandchars=\\\{\}}
% Add ',fontsize=\small' for more characters per line
\usepackage{framed}
\definecolor{shadecolor}{RGB}{241,243,245}
\newenvironment{Shaded}{\begin{snugshade}}{\end{snugshade}}
\newcommand{\AlertTok}[1]{\textcolor[rgb]{0.68,0.00,0.00}{#1}}
\newcommand{\AnnotationTok}[1]{\textcolor[rgb]{0.37,0.37,0.37}{#1}}
\newcommand{\AttributeTok}[1]{\textcolor[rgb]{0.40,0.45,0.13}{#1}}
\newcommand{\BaseNTok}[1]{\textcolor[rgb]{0.68,0.00,0.00}{#1}}
\newcommand{\BuiltInTok}[1]{\textcolor[rgb]{0.00,0.23,0.31}{#1}}
\newcommand{\CharTok}[1]{\textcolor[rgb]{0.13,0.47,0.30}{#1}}
\newcommand{\CommentTok}[1]{\textcolor[rgb]{0.37,0.37,0.37}{#1}}
\newcommand{\CommentVarTok}[1]{\textcolor[rgb]{0.37,0.37,0.37}{\textit{#1}}}
\newcommand{\ConstantTok}[1]{\textcolor[rgb]{0.56,0.35,0.01}{#1}}
\newcommand{\ControlFlowTok}[1]{\textcolor[rgb]{0.00,0.23,0.31}{#1}}
\newcommand{\DataTypeTok}[1]{\textcolor[rgb]{0.68,0.00,0.00}{#1}}
\newcommand{\DecValTok}[1]{\textcolor[rgb]{0.68,0.00,0.00}{#1}}
\newcommand{\DocumentationTok}[1]{\textcolor[rgb]{0.37,0.37,0.37}{\textit{#1}}}
\newcommand{\ErrorTok}[1]{\textcolor[rgb]{0.68,0.00,0.00}{#1}}
\newcommand{\ExtensionTok}[1]{\textcolor[rgb]{0.00,0.23,0.31}{#1}}
\newcommand{\FloatTok}[1]{\textcolor[rgb]{0.68,0.00,0.00}{#1}}
\newcommand{\FunctionTok}[1]{\textcolor[rgb]{0.28,0.35,0.67}{#1}}
\newcommand{\ImportTok}[1]{\textcolor[rgb]{0.00,0.46,0.62}{#1}}
\newcommand{\InformationTok}[1]{\textcolor[rgb]{0.37,0.37,0.37}{#1}}
\newcommand{\KeywordTok}[1]{\textcolor[rgb]{0.00,0.23,0.31}{#1}}
\newcommand{\NormalTok}[1]{\textcolor[rgb]{0.00,0.23,0.31}{#1}}
\newcommand{\OperatorTok}[1]{\textcolor[rgb]{0.37,0.37,0.37}{#1}}
\newcommand{\OtherTok}[1]{\textcolor[rgb]{0.00,0.23,0.31}{#1}}
\newcommand{\PreprocessorTok}[1]{\textcolor[rgb]{0.68,0.00,0.00}{#1}}
\newcommand{\RegionMarkerTok}[1]{\textcolor[rgb]{0.00,0.23,0.31}{#1}}
\newcommand{\SpecialCharTok}[1]{\textcolor[rgb]{0.37,0.37,0.37}{#1}}
\newcommand{\SpecialStringTok}[1]{\textcolor[rgb]{0.13,0.47,0.30}{#1}}
\newcommand{\StringTok}[1]{\textcolor[rgb]{0.13,0.47,0.30}{#1}}
\newcommand{\VariableTok}[1]{\textcolor[rgb]{0.07,0.07,0.07}{#1}}
\newcommand{\VerbatimStringTok}[1]{\textcolor[rgb]{0.13,0.47,0.30}{#1}}
\newcommand{\WarningTok}[1]{\textcolor[rgb]{0.37,0.37,0.37}{\textit{#1}}}

\providecommand{\tightlist}{%
  \setlength{\itemsep}{0pt}\setlength{\parskip}{0pt}}\usepackage{longtable,booktabs,array}
\usepackage{calc} % for calculating minipage widths
% Correct order of tables after \paragraph or \subparagraph
\usepackage{etoolbox}
\makeatletter
\patchcmd\longtable{\par}{\if@noskipsec\mbox{}\fi\par}{}{}
\makeatother
% Allow footnotes in longtable head/foot
\IfFileExists{footnotehyper.sty}{\usepackage{footnotehyper}}{\usepackage{footnote}}
\makesavenoteenv{longtable}
\usepackage{graphicx}
\makeatletter
\def\maxwidth{\ifdim\Gin@nat@width>\linewidth\linewidth\else\Gin@nat@width\fi}
\def\maxheight{\ifdim\Gin@nat@height>\textheight\textheight\else\Gin@nat@height\fi}
\makeatother
% Scale images if necessary, so that they will not overflow the page
% margins by default, and it is still possible to overwrite the defaults
% using explicit options in \includegraphics[width, height, ...]{}
\setkeys{Gin}{width=\maxwidth,height=\maxheight,keepaspectratio}
% Set default figure placement to htbp
\makeatletter
\def\fps@figure{htbp}
\makeatother

\KOMAoption{captions}{tableheading}
\makeatletter
\@ifpackageloaded{tcolorbox}{}{\usepackage[many]{tcolorbox}}
\@ifpackageloaded{fontawesome5}{}{\usepackage{fontawesome5}}
\definecolor{quarto-callout-color}{HTML}{909090}
\definecolor{quarto-callout-note-color}{HTML}{0758E5}
\definecolor{quarto-callout-important-color}{HTML}{CC1914}
\definecolor{quarto-callout-warning-color}{HTML}{EB9113}
\definecolor{quarto-callout-tip-color}{HTML}{00A047}
\definecolor{quarto-callout-caution-color}{HTML}{FC5300}
\definecolor{quarto-callout-color-frame}{HTML}{acacac}
\definecolor{quarto-callout-note-color-frame}{HTML}{4582ec}
\definecolor{quarto-callout-important-color-frame}{HTML}{d9534f}
\definecolor{quarto-callout-warning-color-frame}{HTML}{f0ad4e}
\definecolor{quarto-callout-tip-color-frame}{HTML}{02b875}
\definecolor{quarto-callout-caution-color-frame}{HTML}{fd7e14}
\makeatother
\makeatletter
\makeatother
\makeatletter
\makeatother
\makeatletter
\@ifpackageloaded{caption}{}{\usepackage{caption}}
\AtBeginDocument{%
\ifdefined\contentsname
  \renewcommand*\contentsname{Table of contents}
\else
  \newcommand\contentsname{Table of contents}
\fi
\ifdefined\listfigurename
  \renewcommand*\listfigurename{List of Figures}
\else
  \newcommand\listfigurename{List of Figures}
\fi
\ifdefined\listtablename
  \renewcommand*\listtablename{List of Tables}
\else
  \newcommand\listtablename{List of Tables}
\fi
\ifdefined\figurename
  \renewcommand*\figurename{Figure}
\else
  \newcommand\figurename{Figure}
\fi
\ifdefined\tablename
  \renewcommand*\tablename{Table}
\else
  \newcommand\tablename{Table}
\fi
}
\@ifpackageloaded{float}{}{\usepackage{float}}
\floatstyle{ruled}
\@ifundefined{c@chapter}{\newfloat{codelisting}{h}{lop}}{\newfloat{codelisting}{h}{lop}[chapter]}
\floatname{codelisting}{Listing}
\newcommand*\listoflistings{\listof{codelisting}{List of Listings}}
\makeatother
\makeatletter
\@ifpackageloaded{caption}{}{\usepackage{caption}}
\@ifpackageloaded{subcaption}{}{\usepackage{subcaption}}
\makeatother
\makeatletter
\@ifpackageloaded{tcolorbox}{}{\usepackage[many]{tcolorbox}}
\makeatother
\makeatletter
\@ifundefined{shadecolor}{\definecolor{shadecolor}{rgb}{.97, .97, .97}}
\makeatother
\makeatletter
\makeatother
\makeatletter
\@ifpackageloaded{fontawesome5}{}{\usepackage{fontawesome5}}
\makeatother
\ifLuaTeX
  \usepackage{selnolig}  % disable illegal ligatures
\fi
\IfFileExists{bookmark.sty}{\usepackage{bookmark}}{\usepackage{hyperref}}
\IfFileExists{xurl.sty}{\usepackage{xurl}}{} % add URL line breaks if available
\urlstyle{same} % disable monospaced font for URLs
\hypersetup{
  pdftitle={Bonnes pratiques pour les projets statistiques},
  pdfauthor={Romain Avouac, Thomas Faria, Lino Galiana, Tom Seimandi},
  colorlinks=true,
  linkcolor={blue},
  filecolor={Maroon},
  citecolor={Blue},
  urlcolor={Blue},
  pdfcreator={LaTeX via pandoc}}

\title{Bonnes pratiques pour les projets statistiques}
\usepackage{etoolbox}
\makeatletter
\providecommand{\subtitle}[1]{% add subtitle to \maketitle
  \apptocmd{\@title}{\par {\large #1 \par}}{}{}
}
\makeatother
\subtitle{\textbf{{1ère partie : formation à Git}}}
\author{\href{https://github.com/avouacr}{Romain Avouac},
\href{https://github.com/ThomasFaria}{Thomas Faria},
\href{https://www.linogaliana.fr/}{Lino Galiana},
\href{https://github.com/tomseimandi/}{Tom Seimandi}}
\date{}

\begin{document}
\maketitle
\ifdefined\Shaded\renewenvironment{Shaded}{\begin{tcolorbox}[sharp corners, boxrule=0pt, borderline west={3pt}{0pt}{shadecolor}, breakable, interior hidden, enhanced, frame hidden]}{\end{tcolorbox}}\fi

\hypertarget{introduction}{%
\section{Introduction}\label{introduction}}

\hypertarget{la-notion-de-bonnes-pratiques}{%
\subsection{La notion de bonnes
pratiques}\label{la-notion-de-bonnes-pratiques}}

\begin{itemize}
\tightlist
\item
  {\textbf{Origine}} : communauté des développeurs logiciels
\end{itemize}

. . .

\begin{itemize}
\tightlist
\item
  {\textbf{Constats}} :

  \begin{itemize}
  \tightlist
  \item
    le {\emph{``code est plus souvent lu qu'écrit''}}
    (\href{https://fr.wikipedia.org/wiki/Guido_van_Rossum}{Guido Van
    Rossum})
  \item
    la maintenance d'un code est très coûteuse
  \end{itemize}
\end{itemize}

. . .

\begin{itemize}
\tightlist
\item
  {\textbf{Conséquence}} : ensemble de {\textbf{règles informelles}},
  conventionnellement acceptées comme produisant des logiciels
  {\textbf{fiables}}, {\textbf{évolutifs}} et {\textbf{maintenables}}
\end{itemize}

\hypertarget{pourquoi-sintuxe9resser-aux-bonnes-pratiques}{%
\subsection{Pourquoi s'intéresser aux bonnes pratiques
?}\label{pourquoi-sintuxe9resser-aux-bonnes-pratiques}}

L'activité du statisticien / \emph{datascientist} tend à se rapprocher
de celle du développeur :

. . .

\begin{itemize}
\tightlist
\item
  projets \textbf{intenses en code}
\end{itemize}

. . .

\begin{itemize}
\tightlist
\item
  \textbf{projets collaboratifs} et de grande envergure
\end{itemize}

. . .

\begin{itemize}
\tightlist
\item
  \textbf{complexification} des données et donc des
  \textbf{infrastructures}
\end{itemize}

. . .

\begin{itemize}
\tightlist
\item
  \textbf{déploiement} d'applications pour valoriser les analyses
\end{itemize}

\hypertarget{bonnes-pratiques-et-reproductibilituxe9}{%
\subsection{Bonnes pratiques et
reproductibilité}\label{bonnes-pratiques-et-reproductibilituxe9}}

\begin{figure}

{\centering \includegraphics[width=\textwidth,height=2.08333in]{img/reprospectrum.png}

}

\end{figure}

\textbf{Source} : Peng R., Reproducible Research in Computational
Science, Science (2011)

\begin{itemize}
\item
  Une reproductibilité parfaite est {\textbf{coûteuse}}
\item
  \texttt{Git} est un {\textbf{standard atteignable}} et
  {\textbf{efficient}}
\end{itemize}

\hypertarget{bonnes-pratiques-et-reproductibilituxe9-1}{%
\subsection{Bonnes pratiques et
reproductibilité}\label{bonnes-pratiques-et-reproductibilituxe9-1}}

\begin{figure}

{\centering \includegraphics{img/aea.png}

}

\end{figure}

\begin{itemize}
\tightlist
\item
  Les journaux académiques sont de plus en plus exigeants sur la
  reproductibilité (notamment
  \href{https://aeadataeditor.github.io/}{AEA} et
  \href{https://twitter.com/AeaData}{@AeaData})
\end{itemize}

\hypertarget{le-contruxf4le-de-version-pourquoi-faire}{%
\section{Le contrôle de version : pourquoi faire
?}\label{le-contruxf4le-de-version-pourquoi-faire}}

\hypertarget{archiver-son-code-proprement}{%
\subsection{\texorpdfstring{{1️⃣} Archiver son code
proprement}{1️⃣ Archiver son code proprement}}\label{archiver-son-code-proprement}}

pour en finir avec ça :

\includegraphics{img/fichiers_multiples.png}

\hypertarget{archiver-son-code-proprement-1}{%
\subsection{\texorpdfstring{{1️⃣} Archiver son code
proprement}{1️⃣ Archiver son code proprement}}\label{archiver-son-code-proprement-1}}

ou ça :

\begin{figure}

{\centering \includegraphics{img/finalfinal.png}

}

\end{figure}

\hypertarget{archiver-son-code-proprement-2}{%
\subsection{\texorpdfstring{{1️⃣} Archiver son code
proprement}{1️⃣ Archiver son code proprement}}\label{archiver-son-code-proprement-2}}

ou encore ça :

\begin{Shaded}
\begin{Highlighting}[numbers=left,,]
\NormalTok{prior }\OperatorTok{\textless{}{-}}\NormalTok{ read\_csv(prior\_path)}
\NormalTok{prior }\OperatorTok{\textless{}{-}}\NormalTok{ prior }\OperatorTok{\%\textgreater{}\%}
\NormalTok{    select(}\BuiltInTok{id}\NormalTok{, proba\_inter, proba\_build, proba\_rfl) }\OperatorTok{\%\textgreater{}\%}
\NormalTok{    separate(}\BuiltInTok{id}\NormalTok{, into }\OperatorTok{=}\NormalTok{ c(}\StringTok{\textquotesingle{}nidt\textquotesingle{}}\NormalTok{, }\StringTok{\textquotesingle{}grid\_id\textquotesingle{}}\NormalTok{), sep }\OperatorTok{=} \StringTok{":"}\NormalTok{) }\OperatorTok{\%\textgreater{}\%}
\NormalTok{    group\_by(nidt) }\OperatorTok{\%\textgreater{}\%}
\NormalTok{    mutate(}
\NormalTok{        proba\_build }\OperatorTok{=}\NormalTok{ proba\_build}\OperatorTok{/}\BuiltInTok{sum}\NormalTok{(proba\_build),}
\NormalTok{        proba\_rfl }\OperatorTok{=}\NormalTok{ proba\_rfl}\OperatorTok{/}\BuiltInTok{sum}\NormalTok{(proba\_rfl),}
\NormalTok{        ) }\OperatorTok{\%\textgreater{}\%}
\NormalTok{    unite(col }\OperatorTok{=} \StringTok{"id"}\NormalTok{, nidt, grid\_id, sep }\OperatorTok{=} \StringTok{":"}\NormalTok{)}

\CommentTok{\# Test}
\CommentTok{\# prior\_test \textless{}{-} prior \%\textgreater{}\%}
\CommentTok{\#    mutate(}
\CommentTok{\#        proba\_inter = round(proba\_inter, 4)}
\CommentTok{\#        proba\_build = round(proba\_build, 4)}
\CommentTok{\#        proba\_rfl = round(proba\_rfl, 4)}
\CommentTok{\#    )}

\NormalTok{write\_csv(prior\_round, }\StringTok{"\textasciitilde{}/prior.csv"}\NormalTok{)}
\end{Highlighting}
\end{Shaded}

\hypertarget{archiver-son-code-proprement-3}{%
\subsection{\texorpdfstring{{1️⃣} Archiver son code
proprement}{1️⃣ Archiver son code proprement}}\label{archiver-son-code-proprement-3}}

Pour arriver à ça :

\begin{figure}

{\centering \includegraphics[width=\textwidth,height=4.94792in]{img/timeline.png}

}

\end{figure}

Source :
\href{https://thinkr.fr/travailler-avec-git-via-rstudio-et-versionner-son-code/}{ThinkR}

\hypertarget{voyager-dans-le-temps-de-votre-projet}{%
\subsection{\texorpdfstring{{2️⃣} Voyager dans le temps (de votre
projet)}{2️⃣ Voyager dans le temps (de votre projet)}}\label{voyager-dans-le-temps-de-votre-projet}}

\begin{figure}

{\centering \includegraphics{img/historique_utilitr.png}

}

\end{figure}

\hypertarget{une-collaboration-simplifiuxe9e-et-efficace}{%
\subsection{\texorpdfstring{{3️⃣} Une collaboration simplifiée et
efficace}{3️⃣ Une collaboration simplifiée et efficace}}\label{une-collaboration-simplifiuxe9e-et-efficace}}

Un modèle distribué

\begin{figure}

{\centering \includegraphics[width=\textwidth,height=4.94792in]{img/git_distributed.png}

}

\end{figure}

Source : \href{https://www.specbee.com}{specbee.com}

\hypertarget{une-collaboration-simplifiuxe9e-et-efficace-1}{%
\subsection{\texorpdfstring{{3️⃣} Une collaboration simplifiée et
efficace}{3️⃣ Une collaboration simplifiée et efficace}}\label{une-collaboration-simplifiuxe9e-et-efficace-1}}

Qui permet l'expérimentation en toute sécurité

\begin{figure}

{\centering \includegraphics{img/branches.png}

}

\end{figure}

Source :
\href{https://lutece.paris.fr/support/jsp/site/Portal.jsp?page=wiki\&view=page\&page_name=git\&language=fr}{lutece.paris.fr}

\hypertarget{une-collaboration-simplifiuxe9e-et-efficace-2}{%
\subsection{\texorpdfstring{{3️⃣} Une collaboration simplifiée et
efficace}{3️⃣ Une collaboration simplifiée et efficace}}\label{une-collaboration-simplifiuxe9e-et-efficace-2}}

Avec des outils pour faciliter la collaboration

\begin{figure}

{\centering \includegraphics[width=\textwidth,height=5.20833in]{img/issue.png}

}

\end{figure}

\hypertarget{partager-son-code-uxe0-un-public-large}{%
\subsection{\texorpdfstring{{4️⃣} Partager son code à un public
large}{4️⃣ Partager son code à un public large}}\label{partager-son-code-uxe0-un-public-large}}

\begin{itemize}
\tightlist
\item
  Faciliter la réutilisation
\end{itemize}

\begin{figure}

{\centering \includegraphics[width=\textwidth,height=4.16667in]{img/ines.png}

}

\end{figure}

\hypertarget{partager-son-code-uxe0-un-public-large-1}{%
\subsection{\texorpdfstring{{4️⃣} Partager son code à un public
large}{4️⃣ Partager son code à un public large}}\label{partager-son-code-uxe0-un-public-large-1}}

\begin{itemize}
\tightlist
\item
  Faciliter la réutilisation
\item
  Permettre les contributions extérieures
\end{itemize}

\begin{figure}

{\centering \includegraphics[width=\textwidth,height=3.64583in]{img/utilitr.png}

}

\end{figure}

\hypertarget{partager-son-code-uxe0-un-public-large-2}{%
\subsection{\texorpdfstring{{4️⃣} Partager son code à un public
large}{4️⃣ Partager son code à un public large}}\label{partager-son-code-uxe0-un-public-large-2}}

\begin{itemize}
\tightlist
\item
  Faciliter la réutilisation
\item
  Permettre les contributions extérieures
\item
  Une vitrine pour vous
\end{itemize}

\begin{figure}

{\centering \includegraphics[width=\textwidth,height=3.125in]{img/hadley.png}

}

\end{figure}

\hypertarget{en-ruxe9sumuxe9}{%
\subsection{En résumé}\label{en-ruxe9sumuxe9}}

\begin{itemize}
\tightlist
\item
  {\textbf{Construire}} et {\textbf{naviguer}} à travers
  l'{\textbf{historique}} de son projet
\end{itemize}

. . .

\begin{itemize}
\tightlist
\item
  La {\textbf{collaboration}} rendue {\textbf{simple}} et
  {\textbf{efficace}}
\end{itemize}

. . .

\begin{itemize}
\tightlist
\item
  Améliorer la {\textbf{reproductibilité}} de ses projets
\end{itemize}

. . .

\begin{itemize}
\tightlist
\item
  Améliorer la {\textbf{visibilité}} de ses projets
\end{itemize}

La construction d'un historique du projet permet grandement de diminuer
les risques lors de la maintenance et de la modification du code.

\hypertarget{le-contruxf4le-de-version-avec-git}{%
\section{\texorpdfstring{Le contrôle de version avec
\texttt{Git}}{Le contrôle de version avec Git}}\label{le-contruxf4le-de-version-avec-git}}

\hypertarget{git-est-complexe}{%
\subsection{\texorpdfstring{{⚠️} Git est
complexe}{⚠️ Git est complexe}}\label{git-est-complexe}}

L'utilisation de \texttt{Git} nécessite {\textbf{certaines notions
préalables}}:

\begin{itemize}
\tightlist
\item
  Fonctionnement d'un \texttt{filesystem}
\item
  Connaissance basique du terminal \texttt{Linux}
\item
  Potentiellement, un grand nombre de commandes
\end{itemize}

\begin{figure}

{\centering \includegraphics[width=\textwidth,height=4.16667in]{img/gitfire.png}

}

\caption{Source : \href{https://iulianacosmina.com/?p=19452}{Ici}}

\end{figure}

\hypertarget{git-est-complexe-1}{%
\subsection{\texorpdfstring{{⚠️} Git est
complexe}{⚠️ Git est complexe}}\label{git-est-complexe-1}}

{\textbf{Mais}}

\begin{itemize}
\tightlist
\item
  L'\textbf{usage quotidien} n'implique que {\textbf{quelques
  commandes}}
\item
  {\textbf{Enormément de ressources}} disponibles sur internet
\item
  Des {\textbf{interfaces visuelles}} (ex: \texttt{GitHub\ Desktop},
  \texttt{RStudio}, \texttt{Sublime\ Merge}, \texttt{VS\ Code}) qui
  facilitent l'apprentissage
\item
  Un petit investissement individuel pour de {\textbf{gros gains
  collectifs}}
\end{itemize}

\hypertarget{exemples-dinterfaces-visuelles}{%
\subsection{Exemples d'interfaces
visuelles}\label{exemples-dinterfaces-visuelles}}

\begin{figure}

{\centering \includegraphics{img/interface_github_desk.PNG}

}

\end{figure}

\hypertarget{exemples-dinterfaces-visuelles-1}{%
\subsection{Exemples d'interfaces
visuelles}\label{exemples-dinterfaces-visuelles-1}}

\begin{figure}

{\centering \includegraphics{img/interface_Rstudio.PNG}

}

\end{figure}

\hypertarget{exemples-dinterfaces-visuelles-2}{%
\subsection{Exemples d'interfaces
visuelles}\label{exemples-dinterfaces-visuelles-2}}

\begin{figure}

{\centering \includegraphics{img/interface_commande_prompt.PNG}

}

\end{figure}

\hypertarget{concepts}{%
\subsection{Concepts}\label{concepts}}

\texttt{Git}, \texttt{GitHub}, \texttt{GitLab}, \texttt{GitHub\ Desktop}
\ldots{} quelles différences ?

\begin{itemize}
\tightlist
\item
  \texttt{Git} est un \textbf{logiciel} ;
\item
  Il interprète les lignes de commandes.
\item
  Différentes {\textbf{interfaces graphiques}} permettent d'interagir
  avec Git plus simplement en ``clic-boutton''
  (\texttt{GitHub\ Desktop}, \texttt{RStudio}, \texttt{VS\ Code}\ldots)
\end{itemize}

\hypertarget{concepts-1}{%
\subsection{Concepts}\label{concepts-1}}

\texttt{Git}, \texttt{GitHub}, \texttt{GitLab}, \texttt{GitHub\ Desktop}
\ldots{} quelles différences ?

\begin{itemize}
\tightlist
\item
  \texttt{GitHub} et \texttt{GitLab} sont des \textbf{forges
  logicielles}
\item
  \emph{Forge}: espace d'archivage de code
\item
  Des fonctionalités supplémentaires : \textbf{réseau social du code}
\end{itemize}

\hypertarget{concepts-2}{%
\subsection{Concepts}\label{concepts-2}}

Dépôt local / dépôt distant (\texttt{remote})

\begin{figure}

{\centering \includegraphics{img/localremote.png}

}

\end{figure}

Source :
\href{https://collaboratif-git-formation-insee.netlify.app/des-bases-de-git.html}{Collaborative
work with R}

\hypertarget{concepts-3}{%
\subsection{Concepts}\label{concepts-3}}

\emph{Workflow} (version littéraire) :

\begin{itemize}
\tightlist
\item
  On travaille sur un dépôt local en éditant des fichiers
\item
  On dit à \texttt{Git} que ces fichiers doivent être suivis
  (\texttt{staging\ area})
\item
  On valide les modifications faites en local (\texttt{commit})
\item
  On soumet les modifications (\texttt{push}) après avoir récupéré la
  version collective (\texttt{pull})
\end{itemize}

\hypertarget{concepts-4}{%
\subsection{Concepts}\label{concepts-4}}

\emph{Workflow} (version imagée complète) :

\includegraphics[width=10.41667in,height=5.20833in]{img/workflow_v2.png}

\hypertarget{commandes-essentielles}{%
\subsection{Commandes essentielles}\label{commandes-essentielles}}

\begin{longtable}[]{@{}ll@{}}
\toprule()
Action & Commande \\
\midrule()
\endhead
Cloner un projet & \texttt{git\ clone\ {[}url-to-git-repo{]}} \\
Afficher les changements & \texttt{git\ status} \\
Retrouver l'URL du dépôt distant & \texttt{git\ remote\ -v} \\
\bottomrule()
\end{longtable}

\hypertarget{commandes-essentielles-1}{%
\subsection{Commandes essentielles}\label{commandes-essentielles-1}}

\begin{longtable}[]{@{}
  >{\raggedright\arraybackslash}p{(\columnwidth - 2\tabcolsep) * \real{0.4444}}
  >{\raggedright\arraybackslash}p{(\columnwidth - 2\tabcolsep) * \real{0.5556}}@{}}
\toprule()
\begin{minipage}[b]{\linewidth}\raggedright
Action
\end{minipage} & \begin{minipage}[b]{\linewidth}\raggedright
Commande
\end{minipage} \\
\midrule()
\endhead
Ajouter des changements à l'index de \texttt{Git} & Un seul fichier :
\texttt{git\ add\ \textless{}file-name\textgreater{}} Tous les fichiers
déjà indexés : \texttt{git\ add\ -u} Tous les fichiers {⚠️} :
\texttt{git\ add\ -A} \\
\bottomrule()
\end{longtable}

\begin{tcolorbox}[enhanced jigsaw, colframe=quarto-callout-important-color-frame, rightrule=.15mm, breakable, colbacktitle=quarto-callout-important-color!10!white, coltitle=black, toptitle=1mm, title=\textcolor{quarto-callout-important-color}{\faExclamation}\hspace{0.5em}{Warning}, titlerule=0mm, leftrule=.75mm, colback=white, bottomrule=.15mm, bottomtitle=1mm, toprule=.15mm, opacitybacktitle=0.6, arc=.35mm, left=2mm, opacityback=0]

La méthode \texttt{git\ add\ -A} peut amener à suivre les modifications
de fichiers qui ne devraient pas l'être (par exemple, des données).

Il est recommandé de bien réfléchir avant de l'utiliser (ou d'avoir un
bon \texttt{.gitignore})

\end{tcolorbox}

\hypertarget{commandes-essentielles-2}{%
\subsection{Commandes essentielles}\label{commandes-essentielles-2}}

\begin{longtable}[]{@{}
  >{\raggedright\arraybackslash}p{(\columnwidth - 2\tabcolsep) * \real{0.4444}}
  >{\raggedright\arraybackslash}p{(\columnwidth - 2\tabcolsep) * \real{0.5556}}@{}}
\toprule()
\begin{minipage}[b]{\linewidth}\raggedright
Action
\end{minipage} & \begin{minipage}[b]{\linewidth}\raggedright
Commande
\end{minipage} \\
\midrule()
\endhead
Faire un \texttt{commit} & \texttt{git\ commit\ -m\ "message"} \\
Pousser les changements locaux (branche \texttt{master}) &
\texttt{git\ push\ origin\ master} \\
Récupérer les changements distants (branche \texttt{master}) &
\texttt{git\ pull\ origin\ master} \\
\bottomrule()
\end{longtable}

\hypertarget{quelques-explications-avant-les-applications}{%
\subsection{Quelques explications avant les
Applications}\label{quelques-explications-avant-les-applications}}

Travail avec Git à la Dares : les choix proposés

\begin{itemize}
\tightlist
\item
  Utilisation du GitLab de l'intranet :
  (https://gitlab.intranet.social.gouv.fr/)
\item
  Favorisation de l'interface graphique Github Desktop
\item
  Possibilité d'utiliser Rstudio comme interface graphique
  (cf.~application annexe)
\end{itemize}

\hypertarget{application-0}{%
\subsection{Application 0}\label{application-0}}

\begin{tcolorbox}[enhanced jigsaw, colframe=quarto-callout-tip-color-frame, rightrule=.15mm, breakable, colbacktitle=quarto-callout-tip-color!10!white, coltitle=black, toptitle=1mm, title={Préparation de l'environnement de travail}, titlerule=0mm, leftrule=.75mm, colback=white, bottomrule=.15mm, bottomtitle=1mm, toprule=.15mm, opacitybacktitle=0.6, arc=.35mm, left=2mm, opacityback=0]

\begin{enumerate}
\def\labelenumi{\arabic{enumi}.}
\tightlist
\item
  Se connecter à son compte
  \href{https://gitlab.intranet.social.gouv.fr/}{\texttt{GitLab}}
\item
  Créer un nouveau dépôt public sur \texttt{GitLab}
\item
  Copier le lien HTTPS du projet
\item
  Lancer \texttt{GitHub\ Desktop}
\item
  Cloner le dépôt distant sur votre environnement local :

  \begin{itemize}
  \tightlist
  \item
    \texttt{File} → \texttt{Clone\ repository...} → \texttt{URL}
  \end{itemize}
\item
  Coller le lien du projet et selectionner le chemin où vous souhaitez
  le stocker en local
\item
  Rentrer ses identifiants GitLab si demandés
\end{enumerate}

\end{tcolorbox}

❓ \textbf{Question} : qu'est ce qui différencie le projet cloné d'un
projet quelconque ?

\hypertarget{application-1}{%
\subsection{Application 1}\label{application-1}}

\begin{tcolorbox}[enhanced jigsaw, colframe=quarto-callout-tip-color-frame, rightrule=.15mm, breakable, colbacktitle=quarto-callout-tip-color!10!white, coltitle=black, toptitle=1mm, title={Premiers commits}, titlerule=0mm, leftrule=.75mm, colback=white, bottomrule=.15mm, bottomtitle=1mm, toprule=.15mm, opacitybacktitle=0.6, arc=.35mm, left=2mm, opacityback=0]

\begin{enumerate}
\def\labelenumi{\arabic{enumi}.}
\tightlist
\item
  Créer un dossier 📁 \texttt{scripts}
\item
  Y créer les fichiers \texttt{script1.R} et \texttt{script2.R}, chacun
  contenant quelques commandes \texttt{R} de votre choix
\item
  Ajouter ces fichiers à l'index (\emph{staging area}) de Git en les
  cochant dans l'interface \texttt{GitHub\ Desktop} (ils sont en fait
  cochés par défaut dans l'onglet ``Change'' à gauche)
\item
  Effectuer un \texttt{commit}, auquel on donnera un message descriptif
  pertinent
\item
  Supprimer le fichier \texttt{script1.R} et modifier le contenu du
  fichier \texttt{script2.R}
\item
  Analyser ce qui se passe dans l'interface \texttt{GitHub\ Desktop}
\item
  Effectuer un nouveau \emph{commit} pour ajouter ces modifications à
  l'historique
\end{enumerate}

\end{tcolorbox}

. . .

❓ \textbf{Question} : à ce stade, le dépôt du projet sur
\texttt{GitLab} (\texttt{remote}) a-t-il été modifié ?

\hypertarget{application-2}{%
\subsection{Application 2}\label{application-2}}

\begin{tcolorbox}[enhanced jigsaw, colframe=quarto-callout-tip-color-frame, rightrule=.15mm, breakable, colbacktitle=quarto-callout-tip-color!10!white, coltitle=black, toptitle=1mm, title={Interactions avec le dépôt distant}, titlerule=0mm, leftrule=.75mm, colback=white, bottomrule=.15mm, bottomtitle=1mm, toprule=.15mm, opacitybacktitle=0.6, arc=.35mm, left=2mm, opacityback=0]

\begin{enumerate}
\def\labelenumi{\arabic{enumi}.}
\tightlist
\item
  Effectuer un \texttt{push} pour intégrer les changements locaux au
  projet distant
\item
  Parcourir l'historique du projet sur \texttt{GitLab} (onglet
  \texttt{Repository/commits})

  \begin{enumerate}
  \def\labelenumii{\alph{enumii}.}
  \tightlist
  \item
    Faire apparaître les différences entre deux versions consécutives du
    projet
  \item
    Afficher une version passée du projet
  \end{enumerate}
\item
  Modifier/ajouter le fichier \texttt{README.md} via l'interface
  d'édition web de GitLab
\item
  Faire une modification du projet local, puis
  \texttt{commit}\&\texttt{push}. Que se passe-t-il ?
\item
  Faire un \texttt{fetch} puis un \texttt{pull} pour intégrer les
  changements distants au projet local
\item
  Effectuer à nouveau l'étape de \texttt{push} des changements locaux
\end{enumerate}

\end{tcolorbox}

. . .

❓ \textbf{Question} : comment s'assurer que notre projet local est
toujours iso au projet distant ?

\hypertarget{bonnes-pratiques}{%
\subsection{Bonnes pratiques}\label{bonnes-pratiques}}

\textbf{Que versionne-t-on ?}

\begin{itemize}
\tightlist
\item
  Essentiellement du {\textbf{code source}}
\item
  {\textbf{Pas d'outputs}} (fichiers \texttt{.html}, \texttt{.pdf},
  modèles\ldots)
\item
  {\textbf{Pas de données}}, d'informations locales ou sensibles
\end{itemize}

\begin{tcolorbox}[enhanced jigsaw, colframe=quarto-callout-note-color-frame, rightrule=.15mm, breakable, colbacktitle=quarto-callout-note-color!10!white, coltitle=black, toptitle=1mm, title=\textcolor{quarto-callout-note-color}{\faInfo}\hspace{0.5em}{Note}, titlerule=0mm, leftrule=.75mm, colback=white, bottomrule=.15mm, bottomtitle=1mm, toprule=.15mm, opacitybacktitle=0.6, arc=.35mm, left=2mm, opacityback=0]

Pour définir des règles qui évitent de committer tel ou tel fichier, on
utilise un fichier nommé \textbf{\texttt{.gitignore}}.

Si on mélange du code et des éléments annexes (\emph{output},
données\ldots) dans un même dossier, il {\textbf{faut consacrer du temps
à ce fichier}}.

Le site
\href{https://www.toptal.com/developers/gitignore}{\texttt{gitignore.io}}
peut vous fournir des modèles.

N'hésitez pas à y ajouter des règles conservatrices (par exemple
\texttt{*.csv}), comme cela est expliqué dans
\href{https://www.book.utilitr.org/git.html?q=gitignore\#gitignore}{la
documentation \texttt{utilitR}}.

\end{tcolorbox}

\hypertarget{bonnes-pratiques-1}{%
\subsection{Bonnes pratiques}\label{bonnes-pratiques-1}}

\textbf{Format des commits}

\begin{figure}

\begin{minipage}[t]{0.40\linewidth}

{\centering 

\begin{itemize}
\tightlist
\item
  {\textbf{Fréquence}}

  \begin{itemize}
  \tightlist
  \item
    Aussi souvent que possible
  \item
    Le lot de modifications doit ``faire sens''
  \end{itemize}
\item
  {\textbf{Messages}}

  \begin{itemize}
  \tightlist
  \item
    Courts et informatifs (comme un titre de mail)
  \item
    Décrire \textbf{le pourquoi plutôt que le comment} dans le texte
  \end{itemize}
\end{itemize}

}

\end{minipage}%
%
\begin{minipage}[t]{0.60\linewidth}

{\centering 

\raisebox{-\height}{

\includegraphics{img/titre-commit.png}

}

}

\end{minipage}%

\end{figure}

\hypertarget{application-3}{%
\subsection{Application 3}\label{application-3}}

\begin{tcolorbox}[enhanced jigsaw, colframe=quarto-callout-tip-color-frame, rightrule=.15mm, breakable, colbacktitle=quarto-callout-tip-color!10!white, coltitle=black, toptitle=1mm, title={Le fichier \texttt{.gitignore}}, titlerule=0mm, leftrule=.75mm, colback=white, bottomrule=.15mm, bottomtitle=1mm, toprule=.15mm, opacitybacktitle=0.6, arc=.35mm, left=2mm, opacityback=0]

Le fichier \texttt{.gitignore} a pour objectif d'indiquer à Git les
fichiers ou les types de fichiers dont on ne souhaite pas suivre
(synchroniser) les modifications, et qui doivent donc être ignorés par
git.

\begin{enumerate}
\def\labelenumi{\arabic{enumi}.}
\tightlist
\item
  Créer un fichier ``.gitignore'' à la racine du projet.
\item
  Ajouter le contenu du
  \href{https://github.com/github/gitignore/blob/main/R.gitignore}{.gitignore
  standard R}
\item
  Exclure également les fichiers de type \texttt{*.pdf} et
  \texttt{*.html}
\item
  Créer un dossier \texttt{data} à la racine du projet et à l'intérieur
  de celui-ci un fichier \texttt{data/raw.csv} et ajouter au
  \texttt{.gitignore}: \texttt{data/}
\item
  Effectuer un \emph{commit} pour ajouter le fichier \texttt{.gitignore}
  au projet Git
\item
  Vérifier que tout fonctionne comme attendu
\end{enumerate}

\end{tcolorbox}

. . .

❓ \textbf{Question} : que se passe-t-il lorsque l'on ajoute au
\texttt{.gitignore} des fichiers qui ont déjà été \emph{commit} sur le
projet Git ?

\hypertarget{le-travail-collaboratif-avec-git}{%
\section{Le travail collaboratif avec
Git}\label{le-travail-collaboratif-avec-git}}

\hypertarget{outils-pour-le-travail-collaboratif}{%
\subsection{Outils pour le travail
collaboratif}\label{outils-pour-le-travail-collaboratif}}

\begin{itemize}
\tightlist
\item
  L'éco-système \texttt{Git} {\textbf{facilite} le travail collaboratif}

  \begin{itemize}
  \tightlist
  \item
    \texttt{Git} \faIcon{git-alt} : modèle des {\textbf{branches}}
  \item
    \texttt{GitHub} \faIcon{github} / \texttt{GitLab} \faIcon{gitlab} :
    {\textbf{Issues}, \textbf{Pull Requests}, \textbf{Forks}}
  \end{itemize}
\end{itemize}

. . .

\begin{itemize}
\tightlist
\item
  Ces outils ne remplacent pas une {\textbf{bonne définition de
  l'organisation du travail}} en équipe

  \begin{itemize}
  \tightlist
  \item
    Choix d'un \emph{workflow}
  \item
    Droits d'accès
  \item
    Règles de contribution
  \end{itemize}
\end{itemize}

\hypertarget{application-4}{%
\subsection{Application 4}\label{application-4}}

\begin{tcolorbox}[enhanced jigsaw, colframe=quarto-callout-tip-color-frame, rightrule=.15mm, breakable, colbacktitle=quarto-callout-tip-color!10!white, coltitle=black, toptitle=1mm, title={Synchronisation des dépôts}, titlerule=0mm, leftrule=.75mm, colback=white, bottomrule=.15mm, bottomtitle=1mm, toprule=.15mm, opacitybacktitle=0.6, arc=.35mm, left=2mm, opacityback=0]

\begin{enumerate}
\def\labelenumi{\arabic{enumi}.}
\tightlist
\item
  Se mettre par {\textbf{groupes de 3/4 personnes}}
\item
  L'une des personnes crée un dépôt sur \texttt{GitLab}, ajoute un
  fichier README.md et accorde l'accès ``Master'' aux autres personnes
  du projet (en allant dans ``Settings'' puis ``Members'')
\item
  Créer une copie locale (clone) du projet sur son ordinateur, de la
  même manière que précédemment.
\item
  Créer un fichier
  \texttt{\textless{}votre\_nom\textgreater{}-\textless{}votre\_prenom\textgreater{}.md}.
  Écrire dedans trois phrases de son choix sans ponctuation ni
  majuscules, puis \texttt{commit} et \texttt{push} les modifications
\end{enumerate}

\begin{itemize}
\tightlist
\item
  \textbf{À ce stade, une seule personne (la plus rapide) devrait ne pas
  avoir rencontré de difficulté sur le \texttt{push}.} C'est normal !
  Avant d'accepter une modification, \texttt{Git} vérifie en premier
  lieu la cohérence de la branche locale avec le dépôt distant. Le
  premier ayant fait un \texttt{push} a modifié le dépôt commun ; les
  autres doivent intégrer ces modifications dans leur version locale
  (\texttt{pull}) avant de pouvoir proposer un changement.
\end{itemize}

\begin{enumerate}
\def\labelenumi{\arabic{enumi}.}
\setcounter{enumi}{4}
\tightlist
\item
  Pour ceux dont le \texttt{push} a été refusé, effectuer un
  \texttt{fetch} puis un \texttt{pull} des modifications distantes
\item
  Dans \texttt{GitHub\ Desktop}, afficher l'historique du projet et
  regarder la manière dont ont été intégrées les modifications des
  collaborateurs
\item
  Effectuer à nouveau un \texttt{push} de vos modifications locales
\item
  Les derniers membres du groupe devront refaire les étapes précédentes,
  potentiellement plusieurs fois, pour pouvoir \texttt{push} les
  modifications locales
\end{enumerate}

\end{tcolorbox}

. . .

❓ \textbf{Question} : que se serait-il passé si les différents membres
du groupe avaient effectué leurs modifications sur un seul et même
fichier ?

\hypertarget{application-5}{%
\subsection{Application 5}\label{application-5}}

\begin{tcolorbox}[enhanced jigsaw, colframe=quarto-callout-tip-color-frame, rightrule=.15mm, breakable, colbacktitle=quarto-callout-tip-color!10!white, coltitle=black, toptitle=1mm, title={Résoudre les conflits}, titlerule=0mm, leftrule=.75mm, colback=white, bottomrule=.15mm, bottomtitle=1mm, toprule=.15mm, opacitybacktitle=0.6, arc=.35mm, left=2mm, opacityback=0]

\begin{enumerate}
\def\labelenumi{\arabic{enumi}.}
\tightlist
\item
  On se place dans la même configuration que dans l'application
  précédente
\item
  L'une des personnes modifie le contenu de son fichier
  \texttt{\textless{}votre\_nom\textgreater{}-\textless{}votre\_prenom\textgreater{}.md},
  puis \texttt{commit} et \texttt{push} les modifications
\item
  Sans faire de \texttt{pull} préalable, les autres personnes modifient
  également le contenu du fichier de cette personne, puis
  \texttt{commit} et \texttt{push} les modifications
\end{enumerate}

\begin{itemize}
\tightlist
\item
  Le \texttt{push} est rejeté pour la même raison que dans l'application
  précédente : les dépôts ne sont plus synchronisés, il faut
  \texttt{pull} les changements distants au préalable. Mais
  \textbf{cette fois, le \texttt{pull} est également rejeté} : il y a un
  \textbf{conflit} entre l'historique du projet distant et celui du
  projet local. \texttt{Git} nous indique qu'il faut résoudre le conflit
  avant de pouvoir modifier l'historique du projet.
\end{itemize}

\begin{enumerate}
\def\labelenumi{\arabic{enumi}.}
\setcounter{enumi}{4}
\tightlist
\item
  Pour résoudre le conflit, vous pouvez modifier votre fichier
  manuellement en choisissant la version du fichier que vous souhaitez
  conserver, puis \texttt{commit}/\texttt{push} les modifications
\item
  Comme dans l'application précédente, seul le développeur le plus
  rapide parvient à \texttt{push}. Les autres doivent répéter
  l'opération.
\end{enumerate}

\end{tcolorbox}

. . .

❓ \textbf{Question} : comment limiter au maximum la survenue des
conflits d'historique ?

\hypertarget{le-moduxe8le-des-branches}{%
\subsection{Le modèle des branches}\label{le-moduxe8le-des-branches}}

\includegraphics{img/branches.png}

\hypertarget{exemple-dorganisation-le-github-flow}{%
\subsection{\texorpdfstring{Exemple d'organisation : le \emph{GitHub
flow}}{Exemple d'organisation : le GitHub flow}}\label{exemple-dorganisation-le-github-flow}}

\includegraphics{img/ghflow.png}

Description plus détaillée :
\href{https://docs.github.com/en/get-started/quickstart/github-flow}{ici}

\hypertarget{application-6}{%
\subsection{Application 6}\label{application-6}}

\begin{tcolorbox}[enhanced jigsaw, colframe=quarto-callout-tip-color-frame, rightrule=.15mm, breakable, colbacktitle=quarto-callout-tip-color!10!white, coltitle=black, toptitle=1mm, title={Utilisation de branches}, titlerule=0mm, leftrule=.75mm, colback=white, bottomrule=.15mm, bottomtitle=1mm, toprule=.15mm, opacitybacktitle=0.6, arc=.35mm, left=2mm, opacityback=0]

\begin{enumerate}
\def\labelenumi{\arabic{enumi}.}
\tightlist
\item
  Sur le même dépôt local que pour les applications 4 et 5, récupérer
  les mises à jour de la branche principale du dépôt distant
\item
  Créer une branche intitulée
  \texttt{\textless{}votre\_nom\textgreater{}-\textless{}votre\_prénom\textgreater{}}
  depuis GitHub Desktop en cliquant sur ``Current branch'' puis ``New
  branch''
\item
  Effectuer un \texttt{commit} avec les modifications de votre choix,
  puis pousser les changements sur une nouvelle branche du dépôt distant
\item
  Ouvrir une \texttt{Merge\ Request} dans GitLab pour proposer
  d'intégrer vos changements sur la branche principale du dépôt distant
\end{enumerate}

\end{tcolorbox}

. . .

❓ \textbf{Question} : quelle organisation pour merge dans la branche
principale ?

\hypertarget{application-7}{%
\subsection{Application 7}\label{application-7}}

\begin{tcolorbox}[enhanced jigsaw, colframe=quarto-callout-tip-color-frame, rightrule=.15mm, breakable, colbacktitle=quarto-callout-tip-color!10!white, coltitle=black, toptitle=1mm, title={Contribution à un projet open-source}, titlerule=0mm, leftrule=.75mm, colback=white, bottomrule=.15mm, bottomtitle=1mm, toprule=.15mm, opacitybacktitle=0.6, arc=.35mm, left=2mm, opacityback=0]

\textbf{Comment résoudre un problème mentionné dans une ``issue'' d'un
projet Gitlab ?}

\begin{enumerate}
\def\labelenumi{\arabic{enumi}.}
\tightlist
\item
  Sur \texttt{GitLab}, faire un \texttt{fork} du dépôt
  \href{https://gitlab.intranet.social.gouv.fr/leomoquay/projet_collaboratif_formation}{\texttt{git-exo}}
\item
  Constater la présence d'une ``issue'' à résoudre.
\item
  Créer une branche intitulée
  \texttt{\textless{}votre\_nom\textgreater{}-\textless{}votre\_prénom\textgreater{}}
  à partir de votre ``fork''
\item
  Effectuer un \texttt{commit} avec les modifications permettant de
  remédier au problème, puis pousser les changements sur votre
  \texttt{fork}
\item
  Ouvrir une \texttt{Merge\ Request} pour proposer d'intégrer vos
  changements sur la branche principale du dépôt
  \href{https://github.com/avouacr/git-exo.git}{\texttt{git-exo}} et
  mentionner l'issue (avec \#) à laquelle elle répond.
\end{enumerate}

\end{tcolorbox}

❓ \textbf{Question} : pourquoi, avec un \texttt{fork}, est-il très
important de toujours effectuer une \texttt{Merge\ Request} à partir
d'une branche différente de la branche principale ?

\hypertarget{application-8}{%
\subsection{Application 8}\label{application-8}}

\begin{tcolorbox}[enhanced jigsaw, colframe=quarto-callout-tip-color-frame, rightrule=.15mm, breakable, colbacktitle=quarto-callout-tip-color!10!white, coltitle=black, toptitle=1mm, title={Revenir à des anciennes versions du code}, titlerule=0mm, leftrule=.75mm, colback=white, bottomrule=.15mm, bottomtitle=1mm, toprule=.15mm, opacitybacktitle=0.6, arc=.35mm, left=2mm, opacityback=0]

\textbf{Cette application nécessite un peu de préparation:}

\begin{enumerate}
\def\labelenumi{\arabic{enumi}.}
\tightlist
\item
  Créer un dépôt vide sur le Gitlab et cloner le dépôt sur votre PC
\item
  Créer un fichier ``fichier1'' du type de votre choix puis effectuer
  une commit (commit 1)
\item
  Créer un deuxième fichier, modifier le contenu du fichier 1 et
  effectuer une commit (commit 2)
\item
  Créer un troisième fichier, modifier le contenu du fichier 2 et
  effectuer une commit (commit 3)
\end{enumerate}

\end{tcolorbox}

\begin{tcolorbox}[enhanced jigsaw, colframe=quarto-callout-tip-color-frame, rightrule=.15mm, breakable, colbacktitle=quarto-callout-tip-color!10!white, coltitle=black, toptitle=1mm, title=\textcolor{quarto-callout-tip-color}{\faLightbulb}\hspace{0.5em}{Tip}, titlerule=0mm, leftrule=.75mm, colback=white, bottomrule=.15mm, bottomtitle=1mm, toprule=.15mm, opacitybacktitle=0.6, arc=.35mm, left=2mm, opacityback=0]

Dans la suite de l'application, nous aurons besoin d'identifier les
commits par leur id. Pour trouver l'id d'une commit, on peut utiliser la
commande \texttt{git\ log}.

\end{tcolorbox}

\hypertarget{application-8---suite}{%
\subsection{Application 8 - Suite}\label{application-8---suite}}

\begin{tcolorbox}[enhanced jigsaw, colframe=quarto-callout-tip-color-frame, rightrule=.15mm, breakable, colbacktitle=quarto-callout-tip-color!10!white, coltitle=black, toptitle=1mm, title={Revenir à des anciennes versions du code}, titlerule=0mm, leftrule=.75mm, colback=white, bottomrule=.15mm, bottomtitle=1mm, toprule=.15mm, opacitybacktitle=0.6, arc=.35mm, left=2mm, opacityback=0]

\textbf{Envisageons désormais différents scénarios}

\begin{enumerate}
\def\labelenumi{\arabic{enumi}.}
\tightlist
\item
  Supposons que l'on souhaite annuler la dernière commit \textbf{(avant
  d'avoir push)}:

  \begin{itemize}
  \tightlist
  \item
    Cliquer sur le bouton \texttt{undo} de l'interface GitHubDesktop (en
    bass à gauche du menu Changes) et vérifier que la commit a été
    annulée, mais que le fichier 3 existe toujours.
  \item
    En fait, si l'on souhaite aussi annuler les modifications apportées
    par une commit, il faut ``discard'' les changements apportées par
    cette commit sur chaque fichier (clic droit sur le fichier
    -\textgreater{} discard).
  \end{itemize}
\item
  Supposons maintenant qu'on souhaite revenir à une ancienne version du
  projet (une ancienne commit) pour explorer une nouvelle possibilité de
  développement: on va \textbf{créer une nouvelle branche et se placer
  par exemple à la commit 2}.

  \begin{itemize}
  \tightlist
  \item
    Dans l'historique des commits GithubDesktop, effectuer un clic-droit
    sur la commit 2 et sélectionner ``Create branch from commit''.
    Rentrer le nom de la nouvelle branche et valider l'opération.
  \item
    Vérifier qu'une nouvelle branche a été créée et qu'elle contient le
    projet tel qu'il était après la commit2
  \item
    Revenir à la branche principale en la sélectionnant directement
    depuis le volet de sélection des branches sur Github Desktop.
  \end{itemize}
\item
  Afin de préparer l'application suivante, recréer un fichier3 puis
  commit (commit 3).
\end{enumerate}

\end{tcolorbox}

\hypertarget{application-8---bonus}{%
\subsection{Application 8 - Bonus}\label{application-8---bonus}}

\begin{tcolorbox}[enhanced jigsaw, colframe=quarto-callout-tip-color-frame, rightrule=.15mm, breakable, colbacktitle=quarto-callout-tip-color!10!white, coltitle=black, toptitle=1mm, title={Revenir à des anciennes versions du code avec l'invite de commandes}, titlerule=0mm, leftrule=.75mm, colback=white, bottomrule=.15mm, bottomtitle=1mm, toprule=.15mm, opacitybacktitle=0.6, arc=.35mm, left=2mm, opacityback=0]

\begin{enumerate}
\def\labelenumi{\arabic{enumi}.}
\tightlist
\item
  Supposons que l'on se rende compte que la modification du fichier 1
  (durant la commit 2) ait posé problème et que l'on souhaite
  \textbf{revenir à sa version de la commit 1 avant dernière commit)} :

  \begin{itemize}
  \tightlist
  \item
    Effectuer l'opération
    \texttt{git\ restore\ nom\_fichier1\ -\/-source\ id\_commit1} et
    vérifier que le fichier 1 est de retour à sa première version. Avec
    \texttt{git\ status}, vérifier qu'il n'est pas enregistré dans
    l'index : statut ``\texttt{not\ staged}'').
  \end{itemize}
\item
  Supposons que les dernières commits soient trop nombreuses ou aient
  donné lieu à des erreurs et que l'on souhaite \textbf{effacer ces
  commits de l'historique de git}

  \begin{itemize}
  \item
    L'opération \texttt{git\ reset\ id\_commit\_1} permet de supprimer
    de l'historique git les commits postérieures à la commit 1 sans pour
    autant effacer les modifications du code qu'elles ont entrainées.
    Effectuer l'opération.
  \item
    Vérifier que les commits ont bien disparu de l'historique
    (\texttt{git\ log}). Vérifier que \texttt{git\ status} liste bien
    les fichiers modifiés par les commits 2 et 3, mais que les fichiers
    affectés ne sont pas ajoutés à l'index.
  \item
    Si ces dernières n'étaient pas désirables, on peut désormais
    effectuer un \texttt{git\ restore} pour effacer les modifications
    qu'elles ont apporté au code (on revient donc à la version du code
    datant de la commit 1).
  \item
    En résumé, \texttt{git\ reset} permet de supprimer des commits de
    l'historique, mais pas de revenir directement à une ancienne version
    du code. Si on souhaite que ce soit le cas, il faut ensuite
    effectuer un \texttt{git\ restore}.
  \end{itemize}
\item
  Recréer deux commits, puis cette fois effectuer un \texttt{git\ push}.

  \begin{itemize}
  \tightlist
  \item
    Il est difficile de modifier l'historique du repo distant après un
    \texttt{git\ push}\ldots{} Si l'on souhaite \textbf{annuler les
    modifications apportées par certaines commits}, on peut cependant
    utiliser \texttt{git\ revert}.
  \item
    Effectuer l'opération :
    \texttt{git\ revert\ HEAD\textasciitilde{}2}, puis constater qu'une
    commit a été créée et a permis d'annuler toutes les modifications
    des 2 dernières commits, mais que celles-ci restent présentes dans
    l'historique.
  \item
    Concrètement, il est rare d'avoir besoin de supprimer des commits de
    l'historique du repo distant, qui est justement là pour garder toute
    trace de travail, même infructueux. Il ne faut donc pas avoir peur
    de push régulièrement ! En pratique, il faut surtout faire attention
    à ne pas push par inadvertance des fichiers sensibles.
  \end{itemize}
\end{enumerate}

\end{tcolorbox}

\hypertarget{application-annexe---git-sur-rstudio}{%
\subsection{Application annexe - Git sur
Rstudio}\label{application-annexe---git-sur-rstudio}}

\begin{tcolorbox}[enhanced jigsaw, colframe=quarto-callout-tip-color-frame, rightrule=.15mm, breakable, colbacktitle=quarto-callout-tip-color!10!white, coltitle=black, toptitle=1mm, title={Préparation de l'environnement de travail sur Rstudio}, titlerule=0mm, leftrule=.75mm, colback=white, bottomrule=.15mm, bottomtitle=1mm, toprule=.15mm, opacitybacktitle=0.6, arc=.35mm, left=2mm, opacityback=0]

\begin{enumerate}
\def\labelenumi{\arabic{enumi}.}
\tightlist
\item
  Se connecter à son compte
  \href{https://gitlab.intranet.social.gouv.fr/}{\texttt{GitLab}} et
  créer un nouveau dépôt public sur \texttt{GitLab}
\item
  \href{https://gitlab.intranet.social.gouv.fr/profile/personal_access_tokens}{Générer
  un \emph{token}} (jeton d'accès) d'authentification. Donnez lui un
  nom, une date d'expiration si nécessaire et cochez ``api''. Attention,
  vous ne pourrez afficher ce jeton d'accès qu'une seule fois.
\item
  Lancer \texttt{RStudio}.
\item
  Utiliser \texttt{terminal} (celui de Rstudio par exemple) pour activer
  le stockage des identifiants :

  \begin{itemize}
  \tightlist
  \item
    \texttt{git\ config\ -\/-global\ credential.helper\ store}
  \end{itemize}
\item
  Cloner le dépôt distant sur votre environnement local avec le lien
  HTTPS de votre projet :

  \begin{itemize}
  \item
    \texttt{File} → \texttt{New\ project} → \texttt{Version\ Control} →
    \texttt{Git}. Entrez vos identifiants GitLab ainsi que le
    \texttt{token} en guise de mots de passe.
  \item
    ATTENTION : il se peut que vous rencontriez un problème avant de
    pouvoir vous identifier (en particulier sur machine virtuelle) avec
    le message d'erreur suivant :
  \item
    ``unable to access SSL certificate problem: unable to get local
    issuer certificate''\\
    Pour régler ce problème, allez dans un \texttt{terminal} et entrez
    la commande :
  \item
    \texttt{git\ config\ -\/-global\ http.sslbackend\ schannel})
  \end{itemize}
\item
  Une fois votre projet cloné, vous remarquerez l'apparition d'un onglet
  \texttt{Git} dans la fenêtre en haut à droite de Rstudio. Pour ouvrir
  l'interface graphique Git de Rstudio, cliquez sur Diff dans l'onglet
  Git.
\item
  Vous avez toutes les clefs en main pour faire du Git sur Rstudio, les
  principes sont exactement les mêmes que dans les applications
  précédente hormis quelques exceptions : vous n'avez pas besoin de
  passer par l'étape ``fetch'' avant de pouvoir ``pull'' les
  modifications du dépôt distant et, Rstudio vous crée automatiquement à
  la racine un fichier .Rproj \& un .gitignore au clonage du projet
\end{enumerate}

\end{tcolorbox}

\hypertarget{ressources-utiles}{%
\subsection{Ressources utiles}\label{ressources-utiles}}

\begin{itemize}
\tightlist
\item
  Pour \textbf{aller plus loin}:

  \begin{itemize}
  \tightlist
  \item
    Formation
    \href{https://collaboratif-git-formation-insee.netlify.app/index.html}{Travail
    collaboratif avec \texttt{R}}
  \item
    Cours
    \href{https://ensae-reproductibilite.netlify.app/}{Reproductibilité
    et bonnes pratiques pour les projets de data science} de
    l'\texttt{ENSAE}
  \item
    La \href{https://www.book.utilitr.org/}{documentation
    \texttt{utilitR}} propose plusieurs chapitres sur \texttt{Git}
  \item
    La \href{https://git-scm.com/book/en/v2}{Bible}
  \end{itemize}
\end{itemize}

. . .

\begin{itemize}
\tightlist
\item
  \textbf{Trouver de l'aide}:

  \begin{itemize}
  \tightlist
  \item
    Pour toute question : le salon Tchap
    \href{https://tchap.gouv.fr/\#/room/\#InseeGitGitlablPtu8f1Frns:agent.finances.tchap.gouv.fr}{Insee-Git-Gitlab}
  \item
    Sollicitez vos collègues utilisateurs de Git !
  \end{itemize}
\end{itemize}

\hypertarget{merci-pour-votre-attention}{%
\subsection{Merci pour votre attention
!}\label{merci-pour-votre-attention}}



\end{document}
